% 王相博 — 中文简历（Awesome-CV / 何泰然风格，面试用）
% 编译: xelatex wangxiangbo_zh.tex

\documentclass[10pt,a4paper]{article}

\usepackage[margin=0.52in]{geometry}
\usepackage{fontspec}
\usepackage[UTF8,fontset=none]{ctex}
\usepackage{microtype}
\usepackage{titlesec}
\usepackage{enumitem}
\usepackage{xcolor}
\usepackage{fontawesome5}
\usepackage[hidelinks]{hyperref}

\IfFileExists{fonts/SourceSans3-Regular.ttf}{%
  \setmainfont{SourceSans3}[
    Path = fonts/,
    Extension = .ttf,
    UprightFont = *-Regular,
    BoldFont = *-Semibold,
    Ligatures = TeX,
  ]
}{}

\IfFontExistsTF{Hiragino Sans GB}{%
  \setCJKmainfont{Hiragino Sans GB}
}{%
  \IfFontExistsTF{PingFang SC}{%
    \setCJKmainfont{PingFang SC}
  }{%
    \setCJKmainfont{STHeiti}
  }%
}

\IfFileExists{fonts/academicons.ttf}{%
  \newfontfamily{\AIfont}{academicons}[Path=fonts/, Extension=.ttf]
  \newcommand{\aiGoogleScholar}{{\AIfont\symbol{"E9D4}}}
}{\usepackage{academicons}}

\definecolor{accent}{HTML}{B91C1C}

\pagestyle{empty}
\setlength{\parindent}{0pt}
\setlength{\parskip}{0pt}
\frenchspacing
\emergencystretch=2em

\titleformat{\section}
  {\color{accent}\bfseries\large}
  {}{0pt}{}
  [\color{accent}\vspace{2pt}\titlerule]
\titlespacing*{\section}{0pt}{8pt}{4pt}

\setlist[itemize]{leftmargin=1.15em, itemsep=1.5pt, parsep=0pt, topsep=1.5pt, label=\textbullet}
\setlist[enumerate]{leftmargin=1.7em, itemsep=4pt, parsep=0pt, topsep=2pt}

\newcommand{\me}[1]{\textbf{#1}}
\newcommand{\cvlink}[2]{\href{#1}{#2}}
\newcommand{\icon}[1]{{\color{accent}\footnotesize #1}}
\newcommand{\contact}[3]{\cvlink{#1}{\icon{#2}\,\small #3}}

\newcommand{\entry}[4]{%
  \par\vspace{1pt}%
  \noindent\textbf{#1}\hfill\textbf{#2}\par
  {\small #3\hfill #4}\par
}

\begin{document}

{\fontsize{22pt}{24pt}\selectfont\bfseries 王相博}

\vspace{2pt}
{\small 杭州电子科技大学 \textperiodcentered\ 通信工程 \textperiodcentered\ 本科}

\vspace{5pt}
\contact{mailto:23200410@hdu.edu.cn}{\faEnvelope}{23200410@hdu.edu.cn}
\;{\color{accent}\textbar}\;
\contact{https://raconiy.github.io/}{\faHome}{raconiy.github.io}
\;{\color{accent}\textbar}\;
\contact{https://github.com/raconiy}{\faGithub}{github.com/raconiy}
\;{\color{accent}\textbar}\;
\contact{https://scholar.google.com/citations?user=8E-eDOIAAAAJ}{\aiGoogleScholar}{Google Scholar}

\section{教育背景}

\entry{杭州电子科技大学}{杭州，中国}
{本科，通信工程}{2023 -- 2027}
\begin{itemize}
  \item 专业排名前 10\%。
\end{itemize}

\vspace{3pt}
\entry{悉尼大学}{悉尼，澳大利亚}
{暑期学校}{2024 年夏}
\begin{itemize}
  \item Fundamentals of Deep Learning、Academic English，两门考核均为 Excellent。
\end{itemize}

\section{研究兴趣}

\noindent\textbf{核心方向。}
计算摄影、三维视觉、世界模型——从非理想采集中恢复可用信号，并建立有几何依据的空间表征。

\vspace{3pt}
\noindent\textbf{核心问题。}
如何从真实退化观测（传感器、编解码、稀疏视角）中恢复高保真结构，并将其转化为可预测的空间世界模型？

\section{荣誉}

2024 \quad 悉尼大学暑期学校两门考核 Excellent。\\
2023--2027 \quad 杭州电子科技大学，专业排名前 10\%。

\section{论文}

{\small *共同一作。\textbf{X.\ Wang} 为本人。}

\vspace{3pt}
{\small\noindent\textbf{\color{accent}会议}}
\begin{enumerate}
  \item[\textbf{[C2]}]
        X.\ Feng$^*$, \me{X.\ Wang}$^*$, T.\ Zhong, C.\ Wang, Y.\ Zhao, T.\ Xu, Z.\ Kuang, F.\ Qin, X.\ Yin, and Y.\ Zhu.
        SR3R: Rethinking Super-Resolution 3D Reconstruction With Feed-Forward Gaussian Splatting.
        \textit{CVPR}, 2026. 共同一作.
        \cvlink{https://openaccess.thecvf.com/content/CVPR2026/papers/Feng_SR3R_Rethinking_Super-Resolution_3D_Reconstruction_With_Feed-Forward_Gaussian_Splatting_CVPR_2026_paper.pdf}{[Paper]}
  \item[\textbf{[C1]}]
        \me{X.\ Wang}, W.\ Jiang, J.\ Wang, Y.\ You, S.\ Fang, and F.\ Wen.
        SwitchCodec: Adaptive Residual-Expert Sparse Quantization for High-Fidelity Neural Audio Coding.
        \textit{ICASSP}, 2026. Oral，一作.
        \cvlink{https://arxiv.org/abs/2601.20362}{[Paper]}
\end{enumerate}

\vspace{2pt}
{\small\noindent\textbf{\color{accent}预印本}}
\begin{enumerate}
  \item[\textbf{[P1]}]
        \me{X.\ Wang}, J.\ Zhang, X.\ Yu, L.\ Lei, and D.\ C.\ Zhang.
        HyperClaim: Fine-Grained Cross-Modal Hypergraph Reasoning for Video Misinformation Detection.
        投稿中，2026. 一作.
        \cvlink{https://arxiv.org/abs/2607.28375}{[Paper]}
\end{enumerate}

\section{科研经历}

\entry{RAWLab，香港科技大学（广州）}{广州，中国}
{研究助理，导师：崔梓腾}{2026.07 -- 至今}
\begin{itemize}
  \item 研究方向：计算摄影、真实场景退化下的成像。
\end{itemize}

\vspace{3pt}
\entry{谢菲尔德大学}{谢菲尔德，英国}
{研究助理，导师：Delvin Ce Zhang}{2026.03 -- 至今}
\begin{itemize}
  \item 研究方向：视频虚假信息检测、跨模态超图推理。
  \item 一作 HyperClaim（投稿中）。
\end{itemize}

\vspace{3pt}
\entry{MILab，杭州电子科技大学}{杭州，中国}
{成员，导师：匡振中}{2025 -- 至今}
\begin{itemize}
  \item 研究方向：前馈高斯溅射、超分辨率三维重建。
  \item 共同一作 SR3R（CVPR 2026）。
\end{itemize}

\vspace{3pt}
\entry{IIPL，杭州电子科技大学}{杭州，中国}
{成员，导师：姜文斌、王进}{2024 -- 至今}
\begin{itemize}
  \item 研究方向：神经音频编码、残差专家向量量化。
  \item 一作 SwitchCodec（ICASSP 2026 Oral）。
\end{itemize}

\section{技能}

\noindent\textbf{编程。} Python, MATLAB, \LaTeX, Linux.\\
\textbf{框架。} PyTorch, OpenCV.

\end{document}
